\SourceLeaf{353}%
Theorie für seinen Ansatz charakteristisch ist. Das hat den weiteren Vorteil, dass sich manche Untersuchungen von eigenständigem Interesse ausführlicher entwickeln lassen, als es beim früheren Plan mit einer gleichmäßigen Behandlung der Themen vereinbar gewesen wäre. Diese Untersuchungen sind mit den übrigen eng verwandt, bereiten sie vor und rücken sie in ein helleres Licht. Das gilt auch dann, wenn sie bei den betreffenden Messungen selbst nicht unmittelbar angewandt wurden.

Zu solchen Untersuchungen gehört insbesondere diejenige, die den Gegenstand der vorliegenden \textit{ersten} Abhandlung bildet. Ihr Hauptinhalt ist eine Methode, nach der sich ein Dreiecksnetz auf der Oberfläche eines Rotationsellipsoids ohne jeden Verlust an Genauigkeit so berechnen lässt, als läge es auf einer Kugeloberfläche. Die Methode beruht auf der Lösung eines viel umfassenderen Problems. Der Verfasser gab sie in einer 1822 verfassten Denkschrift an, die Konferenzrat \textsc{Schumacher} im dritten Heft der \textit{Astronomischen Abhandlungen} unter folgendem Titel zum Druck brachte: \textit{Allgemeine Auflösung der Aufgabe, die Theile einer gegebenen Fläche auf einer andern gegebenen Fläche so abzubilden, dass die Abbildung dem Abgebildeten in den kleinsten Theilen ähnlich wird.} Der Verfasser hat diejenigen Abbildungen einer Fläche auf eine andere, die die angegebene Bedingung erfüllen, \textit{konforme Abbildungen} genannt. Das kürzt die Darstellung ab, und als besonders ergiebiges Hilfsmittel für die Rechnungen der höheren Geodäsie verdienen sie einen eigenen Namen. Das sonst vage Adjektiv ist hier also stets in einem genau bestimmten Sinn zu verstehen. Die \textsc{Mercator}-Projektion und die stereografische Projektion sind bekannte Beispiele für konforme Abbildungen der Kugeloberfläche auf die Ebene.

Es braucht kaum eigens erwähnt zu werden, dass die Ähnlichkeit in den kleinsten (unendlich kleinen) Teilen sorgfältig von der Ähnlichkeit in allen endlichen Teilen zu unterscheiden ist. Letztere lässt sich nur in besonderen Fällen erreichen, wenn sich nämlich die erste Fläche auf die zweite selbst oder auf eine ihr ähnliche abwickeln lässt. Wenn die Konformität im Allgemeinen aber \textit{nur} in der Ähnlichkeit der kleinsten Teile besteht, ist das \textit{Vergrößerungsverhältnis}, also das Verhältnis der einander entsprechenden unendlich kleinen Linien auf beiden Flächen, eine Zahl, die %
